\documentclass[12pt]{article}

\usepackage{amsmath, amssymb}
\usepackage{geometry}
\usepackage{setspace}

\geometry{margin=1in}
\linespread{1.7}

\begin{document}

\begin{center}
\textbf{\Large CSE400 -- Lecture 21 \& 22: Inequalities and Tail Bounds}
\end{center}

\vspace{1em}

\section*{1. Introduction and Outline}

\vspace{0.5em}

Topics covered:

\vspace{0.5em}

Tail and motivation of tail bounds

Markov’s Inequality (definition, proof, example)

Chebyshev’s Inequality (definition, proof, example)

Chernoff Bound and Poisson Tail (definition, proof, example)

Jensen’s Inequality (convexity, proof, example)

Case study and system-level understanding

\vspace{1.5em}

\section*{2. Tails of Random Variables}

\vspace{0.5em}

Definition

\vspace{0.5em}

The tail of a random variable is:

\vspace{0.5em}

\[
P(X > x)
\]

\vspace{1em}

Why tails matter

\vspace{0.5em}

Jobs delayed more than 24 hours

Packets dropped due to full buffer

\vspace{1em}

Key Idea

\vspace{0.5em}

Mean $\rightarrow$ average behavior (easy)

Tail $\rightarrow$ rare extreme events (hard)

\vspace{1em}

Why tails are hard to compute

Requires summing many small probabilities

\vspace{2em}

\section*{3. Tail Examples}

\vspace{0.5em}

Example 1 (Binomial Setting)

\vspace{0.5em}

Jobs distributed among servers

\vspace{0.5em}

Let

\vspace{0.5em}

\[
X \sim \text{Binomial}(n, p)
\]

\vspace{1em}

Tail probability:

\vspace{0.5em}

\[
P(X \geq k)
\]

\vspace{0.5em}

\[
= \sum_{i=k}^{n}
\]

\vspace{0.5em}

\[
\binom{n}{i}
\]

\vspace{0.5em}

\[
p^i (1-p)^{n-i}
\]

\vspace{2em}

Example 2 (Poisson Setting)

\vspace{0.5em}

Jobs arrive with rate

\vspace{0.5em}

\[
X \sim \text{Poisson}(\lambda)
\]

\vspace{1em}

Tail probability:

\vspace{0.5em}

\[
P(X \geq k)
\]

\vspace{0.5em}

\[
= \sum_{i=k}^{\infty}
\]

\vspace{0.5em}

\[
e^{-\lambda}
\]

\vspace{0.5em}

\[
\frac{\lambda^i}{i!}
\]

\vspace{2em}

\section*{4. Why Tail Bounds?}

\vspace{0.5em}

Exact computation:

\vspace{0.5em}

These are complicated long sums

\vspace{1em}

Idea

\vspace{0.5em}

Instead of exact value $\rightarrow$ find an upper bound

Easier to compute

Still provides guarantee

\vspace{2em}

\section*{5. Running Example}

\vspace{0.5em}

Flip a fair coin $n$ times

\vspace{0.5em}

\[
X \sim \text{Binomial}(n, 1/2)
\]

\vspace{0.5em}

\[
E[X]
\]

\vspace{0.5em}

\[
= \frac{n}{2}
\]

\vspace{1em}

Question:

\vspace{0.5em}

\[
P(X \geq \frac{4}{3}n)
\]

\vspace{2em}

\section*{6. Markov’s Inequality}

\vspace{0.5em}

Key Idea

Uses only the mean

Large values are rare

\vspace{1em}

Markov Inequality

\vspace{0.5em}

\[
P(X \geq a)
\]

\vspace{0.5em}

\[
\leq \frac{E[X]}{a}
\]

\vspace{1em}

Proof Idea

\vspace{0.5em}

\[
E[X]
\]

\vspace{0.5em}

\[
= \sum_{x=0}^{\infty} x p_X(x)
\]

\vspace{0.5em}

\[
\geq \sum_{x=a}^{\infty} x p_X(x)
\]

\vspace{0.5em}

\[
\geq a \sum_{x=a}^{\infty} p_X(x)
\]

\vspace{0.5em}

\[
= a \cdot P(X \geq a)
\]

\vspace{1em}

Thus:

\vspace{0.5em}

\[
P(X \geq a)
\]

\vspace{0.5em}

\[
\leq \frac{E[X]}{a}
\]

\vspace{2em}

\section*{7. Chebyshev’s Inequality}

\vspace{0.5em}

Key Idea

Uses mean and variance

Controls deviation from mean

\vspace{1em}

\[
P(|X - \mu| \geq a)
\]

\vspace{0.5em}

\[
\leq \frac{Var(X)}{a^2}
\]

\vspace{1em}

Proof Idea

\vspace{0.5em}

\[
(X - \mu)^2
\]

\vspace{0.5em}

\[
P((X - \mu)^2 \geq a^2)
\]

\vspace{0.5em}

\[
\leq \frac{E[(X - \mu)^2]}{a^2}
\]

\vspace{0.5em}

\[
= \frac{Var(X)}{a^2}
\]

\vspace{2em}

\section*{8. Chernoff Bound}

\vspace{0.5em}

\[
E[e^{tX}]
\]

\vspace{0.5em}

\[
= e^{\lambda(e^t - 1)}
\]

\vspace{1em}

\[
P(X \geq a)
\]

\vspace{0.5em}

\[
\leq \min_{t>0}
\]

\vspace{0.5em}

\[
e^{\lambda(e^t -1) - ta}
\]

\vspace{2em}

\section*{9. Jensen’s Inequality}

\vspace{0.5em}

\[
E[X^2]
\]

\vspace{0.5em}

\[
\geq (E[X])^2
\]

\vspace{1em}

\[
g(\alpha x_1 + (1-\alpha)x_2)
\]

\vspace{0.5em}

\[
\leq \alpha g(x_1) + (1-\alpha) g(x_2)
\]

\vspace{2em}

\section*{10. Case Study}

\vspace{0.5em}

Demand: $X$

Capacity: $C$

\vspace{0.5em}

\[
P(X \geq C)
\]

\vspace{0.5em}

\[
\leq 0.01
\]

\vspace{1em}

\[
P(X \geq C) \leq \frac{\mu}{C}
\]

\vspace{0.5em}

\[
C = 1000000
\]

\vspace{1em}

\[
C = 30000
\]

\end{document}